\section{The model structure and the Quillen equivalence}
\label{sec:transfer}

Fix $n\in\{1,2,\ldots,\infty\}$.
Corollary~\ref{cor:horn-acyclic} supplies the acyclic horn attachments
needed for the model axioms; Corollary~\ref{cor:attachment-invariance}
supplies the invariance under boundary attachment needed for left
properness. We extend these statements through transfinite compositions,
then prove the Quillen equivalence by checking the unit on arbitrary
simplicial sets. Write
\[
 L_n=c_n\Sd^2,\qquad R_n=\Ex^2\Street,
 \qquad L_n\dashv R_n.
\]
We use the Kan--Quillen model structure on $\sSet$, with weak
equivalences $W_{\mathrm{KQ}}$, fibrations
$\mathrm{Fib}_{\mathrm{KQ}}$, and generating sets
\begin{align*}
 I&=\{\partial\Delta[m]\hookrightarrow\Delta[m]\mid m\geq0\},\\
 J&=\{\Lambda^k[m]\hookrightarrow\Delta[m]\mid m\geq1,
                    \ 0\leq k\leq m\}.
\end{align*}
In particular, $\partial\Delta[0]=\varnothing$.

We first record the classical simplicial facts needed to identify the
weak equivalences.  The last-vertex transformation
$\ell:\Sd\longrightarrow\id$ is an objectwise weak equivalence
\cite[Theorem~4.5]{AraMaltsiniotis2014}.  Its mate
$j:\id\longrightarrow\Ex$ is also an objectwise weak equivalence
\cite[Theorem~III.4.6]{GoerssJardine1999}.  The latter assertion is
a theorem about subdivision and extension, and is not a formal
consequence of being a mate.  Both statements apply to every
simplicial set independently of the model structure being constructed.
For a map $f:K\to K'$, naturality gives
\[
 j_{K'}f=(\Ex f)j_K.
\]
Two-out-of-three therefore proves that $\Ex$ preserves and reflects
weak equivalences.  Similarly, naturality of $\ell$ shows that $\Sd$
preserves and reflects weak equivalences.  In particular,
\begin{equation}\label{eq:model-classes}
 \begin{split}
 W_n&=\Street^{-1}(W_{\mathrm{KQ}})
        =R_n^{-1}(W_{\mathrm{KQ}}),\\
 F_n&=R_n^{-1}(\mathrm{Fib}_{\mathrm{KQ}}),\qquad
 C_n={}^{\perp}(F_n\cap W_n).
 \end{split}
\end{equation}
Here ${}^{\perp}$ denotes the left lifting property.  The class $W_n$
has two-out-of-three and is closed under retracts.  By
Lemma~\ref{lem:filtered}, it is also closed under filtered colimits of
arrows, without any injectivity requirement on their transition maps.

For a set $\mathcal S$ of maps, write $\mathcal S\text{-inj}$ for the
maps with the right lifting property against $\mathcal S$, and
$\mathcal S\text{-cof}={}^{\perp}(\mathcal S\text{-inj})$.
A relative $\mathcal S$-cell map is a transfinite composite of
pushouts of coproducts of members of $\mathcal S$.

\begin{lemma}[Transfinite horn attachments]\label{lem:transfinite-horns}
Every relative $L_nJ$-cell map belongs to $W_n$.
\end{lemma}

\begin{proof}
Corollary~\ref{cor:horn-acyclic} proves the assertion for one
attachment, with an arbitrary attaching target.  A set of simultaneous
attachments may be well ordered and attached one at a time: at each
stage use the given map from its source into the original target,
followed by the map to the current stage.  The universal property of
the resulting colimit is the universal property of the pushout of
the coproduct.  Thus it suffices to consider a continuous ordinal
sequence $(X_\alpha)_{\alpha\leq\lambda}$ in which each successor
map is a single such attachment.

By transfinite induction, $X_0\to X_\alpha$ belongs to $W_n$ for
every $\alpha$.  The assertion at zero is the identity, and the
successor step is composition with the weak equivalence supplied by
Corollary~\ref{cor:horn-acyclic}.  If $\beta$ is a nonzero limit
ordinal, the ordinal of stages $\alpha<\beta$ is filtered.  Apply
Lemma~\ref{lem:filtered} to the diagram of arrows
$X_0\to X_\alpha$.  Its colimit is
$X_0\to\colim_{\alpha<\beta}X_\alpha=X_\beta$, because the
colimit of the constant source diagram is $X_0$.  This proves the
limit step and the lemma.
\end{proof}

\begin{proposition}[The transferred model structure]\label{prop:model-structure}
The classes $(C_n,W_n,F_n)$ in
\eqref{eq:model-classes} form a combinatorial model structure on
$\nCat$, cofibrantly generated by $(L_nI,L_nJ)$.
\end{proposition}

\begin{proof}
The underlying category is locally presentable, and hence complete
and cocomplete, by Lemma~\ref{lem:filtered}.  We describe the small
object argument to make clear that the factorizations have no
restriction on their sources.  For either set
$\mathcal S=L_nI$ or $\mathcal S=L_nJ$, choose an infinite regular
cardinal $\kappa$ such that every domain in $\mathcal S$ is
$\kappa$-presentable.  Such a common bound exists because the
category is locally presentable and $\mathcal S$ is a set.

Given any map $X\to Y$, start with $X_0=X$.  At stage $\alpha+1$
attach to $X_\alpha$ one copy of the codomain of each member of
$\mathcal S$ for every commutative square from that member to
$X_\alpha\to Y$.  This is a set of squares, by local smallness.
Use colimits at limit stages.  The resulting factorization is
\[
 X\xrightarrow{i}X_\kappa\xrightarrow{p}Y.
\]
Here $i$ is a relative $\mathcal S$-cell map and
$p\in\mathcal S\text{-inj}$.  Indeed a map from a generating
domain to $X_\kappa$ factors through some $X_\alpha$, since the
ordinal $\kappa$ is $\kappa$-filtered.  Its lifting square against
$p$ was therefore attached at stage $\alpha+1$, which supplies the
required lift.  The construction may be made functorial by using
all squares at each stage.

Pushouts, transfinite composites, and retracts of maps with a fixed
left lifting property retain that property.  Consequently every
relative $\mathcal S$-cell map lies in $\mathcal S\text{-cof}$.
Conversely, if $u:A\to B$ belongs to $\mathcal S\text{-cof}$,
factor it as $A\xrightarrow{i}Z\xrightarrow{p}B$ by the construction
above.  The square with top map $i$, left map $u$, right map $p$,
and bottom map $\id_B$ admits a diagonal $s:B\to Z$.  The equations
$su=i$ and $ps=\id_B$ exhibit $u$ as a retract of $i$ in the arrow
category, with the source fixed at $A$.  Thus the construction gives
the weak factorization system
\[
 (\mathcal S\text{-cof},\mathcal S\text{-inj}),
\]
and its left class consists exactly of retracts of relative
$\mathcal S$-cell maps.

The adjunction identifies lifting squares and their fillers.  The
Kan--Quillen lifting characterizations therefore give exactly
\begin{align}
 (L_nI)\text{-inj}
   &=R_n^{-1}(I\text{-inj})=F_n\cap W_n,
       \label{eq:boundary-injectives}\\
 (L_nJ)\text{-inj}
   &=R_n^{-1}(J\text{-inj})=F_n.
       \label{eq:horn-injectives}
\end{align}
In the first equality of classes on the right of
\eqref{eq:boundary-injectives}, we used the equality of the two
definitions of $W_n$ proved above.  The first weak factorization
system is consequently $(C_n,F_n\cap W_n)$.

Let $\mathcal D=(L_nJ)\text{-cof}$.  As $F_n\cap W_n\subseteq F_n$,
one has $\mathcal D\subseteq C_n$.  Lemma~\ref{lem:transfinite-horns} and
retract closure give $\mathcal D\subseteq W_n$ as well.  To prove the
reverse inclusion, take $u:A\to B$ in $C_n\cap W_n$ and factor it
by the second small object argument as
$A\xrightarrow{i}Z\xrightarrow{p}B$, with $i\in\mathcal D$ and
$p\in F_n$.  Since $i$ and $u=pi$ belong to $W_n$,
two-out-of-three gives $p\in F_n\cap W_n$.  The lifting property
of $u\in C_n$ now supplies $s:B\to Z$ with $su=i$ and
$ps=\id_B$.  The same fixed-source retract diagram puts $u$ in
$\mathcal D$.  We have proved
\[
 \mathcal D=C_n\cap W_n.
\]
Thus the second weak factorization system is $(C_n\cap W_n,F_n)$.
Together with two-out-of-three and retract closure for $W_n$, these
are all the model axioms.  Local presentability and the displayed
generating sets make the model structure combinatorial.
\end{proof}

\begin{proposition}[Properness]\label{prop:properness}
The model structure of Proposition~\ref{prop:model-structure} is
left and right proper.
\end{proposition}

\begin{proof}
For right properness, let $p:Y\to Z$ belong to $F_n$ and let
$w:Z'\to Z$ belong to $W_n$.  The right adjoint $R_n$ carries the
pullback $Y\times_Z Z'$ to the corresponding simplicial pullback.
Its map to $R_nY$ is a weak equivalence by right properness of the
Kan--Quillen model structure, since $R_np$ is a Kan fibration and
$R_nw$ is a weak equivalence.  The definition of $W_n$ then gives
$Y\times_Z Z'\to Y$ in $W_n$.

For left properness, first let $i:X\to Y$ be a relative
$L_nI$-cell map and let $w:X\to X'$ belong to $W_n$.  As in
Lemma~\ref{lem:transfinite-horns}, write $i$ as a continuous
sequence of individual boundary attachments $X_\alpha$ starting
at $X$.  Put $X'_\alpha=X'\amalg_X X_\alpha$.
We prove by induction that $X_\alpha\to X'_\alpha$ belongs to
$W_n$.  This is $w$ at stage zero.  At a successor stage, the
second attachment has the first attaching map followed by
$X_\alpha\to X'_\alpha$, so
Corollary~\ref{cor:attachment-invariance} proves the induction
step.  At a nonzero limit ordinal, both sequences are their
colimits, and Lemma~\ref{lem:filtered} gives the assertion.  Hence
$Y\to X'\amalg_XY$ belongs to $W_n$.

For a general $i:X\to Y$ in $C_n$, the first small object
factorization and its lifting retract, used in
Proposition~\ref{prop:model-structure}, express $i$ as a retract
of a relative $L_nI$-cell map $j:X\to Z$ with the same source.
Explicitly there are $s:Y\to Z$ and $p:Z\to Y$ with
$si=j$, $pj=i$, and $ps=\id_Y$.  Cobase change along $w$ gives
the corresponding maps $s'$ and $p'$ between
$X'\amalg_XY$ and $X'\amalg_XZ$, with $p's'=\id$.
Thus the arrow $Y\to X'\amalg_XY$ is a retract of
$Z\to X'\amalg_XZ$.  The latter is in $W_n$ by the preceding
paragraph, and so is the former.  Both properness arguments apply
to arbitrary objects; no cofibrancy or fibrancy condition on the
base objects was used.
\end{proof}

Let $\eta_T:T\to\Street c_nT$ denote the unit of
$c_n\dashv\Street$.  The poset unit of
Proposition~\ref{prop:poset-unit} is the case $T=\OrdN E$.
The next proposition extends the required comparison to all
simplicial sets after two subdivisions.

\begin{proposition}[The unit after two subdivisions]\label{prop:subdivided-unit}
For every simplicial set $K$, the canonical map
\begin{equation}\label{eq:subdivided-unit}
 \theta_K=\eta_{\Sd^2K}:\Sd^2K\longrightarrow\Street L_nK
\end{equation}
is a weak equivalence.
\end{proposition}

\begin{proof}
For $K=\Delta[m]$ and $K=\partial\Delta[m]$,
Lemma~\ref{lem:subdivision-posets} identifies $\Sd^2K$ with the
ordinary nerve of a poset.  The assertion is therefore
Proposition~\ref{prop:poset-unit}.  It also holds for the empty
simplicial set: $c_n\varnothing$ is the empty strict category,
whose Street nerve is empty.

Every simplicial set is obtained from the empty set by a
continuous transfinite sequence of boundary attachments.  One
construction orders its nondegenerate simplices by dimension and
well orders those in each dimension.  The boundary of a simplex
then maps into the stage already constructed.  That attaching
map may identify vertices or faces, or send faces to degenerate
simplices.  The construction still recovers $K$, by the unique
expression of each simplex as a degeneracy of a nondegenerate
simplex.

Consider one successor step
\[
 K_+=K_0\amalg_{\partial\Delta[m]}\Delta[m]
\]
and assume the assertion for $K_0$.  Write
\[
 V=\Sd^2K_0,\qquad A=\Sd^2\partial\Delta[m],
 \qquad B=\Sd^2\Delta[m].
\]
Here $A,B,V$ are simplicial sets; the representing posets for $A,B$
will be denoted by $E,F$ below. Since subdivision is a left adjoint,
$\Sd^2K_+=V\amalg_A B$.  Naturality of $\eta$ gives a map of
spans and hence the map
\begin{equation}\label{eq:unit-gluing-map}
 q:V\amalg_A B\longrightarrow
       \Street c_nV\amalg_{\Street c_nA}\Street c_nB.
\end{equation}
The three maps on the objects of these spans are weak
equivalences: on $V$ by induction and on $A,B$ by the simplex and
boundary cases.  The leg $A\to B$ is a simplicial monomorphism.
Writing $A=\OrdN E$ and $B=\OrdN F$ as in
Lemma~\ref{lem:subdivision-posets}, the inclusion $E\subseteq F$
is a sieve; Lemma~\ref{lem:poset-pieces} makes
$c_nA\to c_nB$ a sieve inclusion.  Consequently
$\Street c_nA\to\Street c_nB$ is also a monomorphism.
Lemma~\ref{lem:simplicial-gluing}, applied with these two
specified monomorphic legs, proves that $q$ is a weak equivalence.
It requires no monomorphism condition on $A\to V$.

Apply Proposition~\ref{prop:attachment-comparison} to the
arbitrary attaching functor $c_nA\to c_nV$.  It says that
\begin{equation}\label{eq:unit-attachment-comparison}
 t:\Street c_nV\amalg_{\Street c_nA}\Street c_nB
       \longrightarrow
       \Street(c_nV\amalg_{c_nA}c_nB)
\end{equation}
is a weak equivalence.  Because $c_n$ preserves colimits, the
target is $\Street c_n(V\amalg_A B)$.
The composite $tq$ is exactly $\eta_{V\amalg_A B}$: on each of
the two pushout legs $V$ and $B$ this follows from naturality of
$\eta$, and equality on those legs determines the map from the
pushout.  Thus the actual canonical map
$\theta_{K_+}$ is a weak equivalence.  This also explains why
noninjective boundary maps cause no change to the argument.

At a nonzero limit ordinal, $\Sd^2$ and $L_n$ preserve the
colimit as left adjoints, while $\Street$ preserves the filtered
colimit by Lemma~\ref{lem:filtered}.  Therefore the unit at that
stage is the filtered colimit of the previous units; it is a weak
equivalence by the same lemma.  Transfinite induction proves
\eqref{eq:subdivided-unit} for every $K$.  The poset descriptions
used in this proof concern the simplex and boundary pieces only.
\end{proof}

\begin{theorem}[Quillen equivalence]\label{thm:quillen-equivalence}
The adjunction
\[
 L_n:\sSet_{\mathrm{KQ}}\rightleftarrows
       (\nCat,C_n,W_n,F_n):R_n
\]
is a Quillen equivalence.
\end{theorem}

\begin{proof}
Equations~\eqref{eq:boundary-injectives} and
\eqref{eq:horn-injectives} show that $R_n$ preserves fibrations
and trivial fibrations, so this is a Quillen adjunction.

Set $S=\Sd$ and $T=\Ex$, and let $a_K:K\to TSK$ be the unit of
$S\dashv T$. The mate identity is
\[
 j_K=T(\ell_K)a_K:K\longrightarrow TK.
\]
The independently cited theorem above makes $j_K$ a weak equivalence.
Since $\ell_K$ is a weak equivalence and $T$ preserves weak
equivalences, $T(\ell_K):TSK\to TK$ is also one. Hence $a_K$ is a
weak equivalence for every simplicial set $K$, by two-out-of-three.
In particular, $a_{SK}:SK\to TS^2K$ is a weak equivalence. The unit
of $S^2\dashv T^2$ is
\[
 b_K=T(a_{SK})a_K:K\longrightarrow T^2S^2K,
 \qquad T(a_{SK}):TSK\longrightarrow T^2S^2K.
\]
Thus $b_K$ is a weak equivalence because $T$ preserves weak equivalences.
Composition is written with the rightmost map applied first.

The unit of the composite adjunction $L_n\dashv R_n$ is
\[
 u_K=T^2(\theta_K)b_K:
       K\longrightarrow T^2\Street c_nS^2K=R_nL_nK.
\]
Proposition~\ref{prop:subdivided-unit}, preservation of weak
equivalences by $T$, and the preceding calculation show that
$u_K$ is a weak equivalence for every $K$.
For any strict functor $f:L_nK\to Y$, its simplicial adjunct is
\[
 f^\sharp=R_n(f)u_K:K\longrightarrow R_nY.
\]
Using the definition of $W_n$ and two-out-of-three now gives
\[
 f\in W_n
 \quad\Longleftrightarrow\quad
 R_n(f)\in W_{\mathrm{KQ}}
 \quad\Longleftrightarrow\quad
 f^\sharp\in W_{\mathrm{KQ}}.
\]
This equivalence holds for every $K$ and $Y$.  In particular it
holds for cofibrant $K$ and fibrant $Y$, which is the
Quillen-equivalence criterion.  Together with
Propositions~\ref{prop:model-structure} and~\ref{prop:properness},
this proves all the asserted model-categorical conclusions for
each finite $n\geq1$ and for $n=\infty$.
\end{proof}
